According to variation theory \autocite[See][]{NCOL}, learning an educational
objective requires the learner to discern all critical aspects of that
objective, and discernment in turn requires experiencing a pattern of variation
in the dimension of each critical aspect.
In a companion paper we analyse misconceptions in introductory programming
through this lens \autocite{VTProgMisconceptions}.
Marton also uses variation theory to account for learning \emph{without} a
teacher: in scientific discovery, and in the deep approach to learning, the
learner introduces the appropriate patterns of variation for themselves
\autocite{NCOL}.

When a programmer debugs, the bug is an unknown: the programmer's understanding
of the program differs from the program's actual behaviour, so there is a
critical aspect the programmer has not yet discerned.
Debugging is thus learning, and successful debugging should require the
debugger to generate the patterns of variation---contrast, generalisation and
fusion---for themselves.
We hypothesise that debugging skill and the deep approach to learning are
expressions of the same underlying ability: that of introducing the appropriate
patterns of variation for oneself.

We outline a study in an introductory programming course where students' every
edit--run cycle is recorded with the \texttt{learnlog} tool \autocite{learnlog}.
We measure the students' approaches to learning with the R-SPQ-2F questionnaire
\autocite{RSPQ2F} and look for the patterns of variation in the recorded
debugging episodes, relating the two to each other and to debugging success.
